\scanpage{97}
\margin{-5-}

for $L$, but there does not seem to be any simple way of relating 1.4 and 1.5 to their analogues for $L$.

We now introduce a perfect subfield $k$ of $K$, although it appears from recent results of A. Grothendieck on semisimple groups over arbitrary fields that the assumption of perfectness is unnecessary for most of what follows.

1.6 Theorem. \uline{Let} $G$ \uline{be over} $k$, \uline{and assume either that} $G$ \uline{splits over} $k$ \uline{or that} $G$ \uline{contains a Borel subgroup over} $k$ \uline{but no component of type} $A_n$ ($n$ \uline{even}). \uline{Then the set} $N$ \uline{of} 1.4 \uline{can be constructed over} $k$ (\uline{by appropriate choice of} $T$, $\sigma_i$, \uline{etc.}).

Together with 1.4 this implies that if $G$ is simply connected in 1.6 the natural map from the set of regular elements over $k$ to the set of regular classes over $k$ is surjective. For a group of type $A_n$ ($n$ even) we have a substitute (see 9.7) for 1.6 which enables us to show:

1.7 Theorem. \uline{Assume that} $G$ \uline{is simply connected and over} $k$ \uline{and that} $G$ \uline{contains a Borel subgroup over} $k$. \uline{Then the natural map from the set of semisimple elements over} $k$ \uline{to the set of semisimple classes over} $k$ \uline{is surjective. In other words, each semisimple class over} $k$ \uline{contains an element over} $k$.

Theorems 1.6 and 1.7 are proved in \S9 where it is also shown (see 9.1 and 9.10) that the assumption that $G$ contains a Borel subgroup over $k$ is essential.

\mspage{96}
